% 表紙
\begin{titlepage}
\begin{center}
\vspace*{30mm}

{\color{accentblue}\rule{\textwidth}{2pt}}
\vspace{8mm}

{\Huge\bfseries\color{headerblue} \doctitle}

\vspace{6mm}

{\Large\color{accentblue} \docsubtitle}

\vspace{4mm}
{\color{accentblue}\rule{\textwidth}{2pt}}

\vspace{20mm}

\begin{tabular}{ll}
    \textbf{バージョン} & \docversion \\[4pt]
    \textbf{作成日} & \docdate \\[4pt]
    \textbf{著者} & \docauthor \\
\end{tabular}

\vspace{30mm}

\begin{tcolorbox}[
  enhanced,
  colback=accentblue!5,
  colframe=accentblue,
  arc=2mm,
  width=14cm,
  breakable,
]
    \centering
    本書は、ESP32を題材にした組み込みシステム向け\\
    「モジュール設計パターン」をQ\&A形式で解説する教科書です。\\[4pt]
    クラスの基礎から3フェーズ実行モデルまで、\\
    段階的に学べる構成になっています。
\end{tcolorbox}

\vspace{10mm}

\begin{tcolorbox}[
  enhanced,
  colback=noteyellow!5,
  colframe=noteyellow!70,
  arc=2mm,
  width=14cm,
  breakable,
]
    \centering
    {\bfseries 対象読者と読み方}\\[4pt]
    「C++のクラスはなんとなく知ってるけど、\\部品化の経験がない」方を主な対象としています。\\[4pt]
    第1章「前提知識」ではクラスの基礎を解説しています。\\[2pt]
    構造体・クラス・コンストラクタを既にご存知の方は、\\第1章を飛ばして第2章から読み始めてください。
\end{tcolorbox}

\vfill
\end{center}
\end{titlepage}
